% ---------------------------------------------------------------------
\subsection{Individuazione e caratterizzazione degli utenti del sistema}
\label{sec:caratterizzazione-utenti}
% ---------------------------------------------------------------------

Il sistema \productname{} è progettato per supportare tre distinte tipologie di utenti, ciascuna con obiettivi, responsabilità e permessi specifici. Questa sezione fornisce una caratterizzazione dettagliata di ogni profilo utente.

% ---------------------------------------------------------------------
\subsubsection{Utente}
\label{subsec:utente}
% ---------------------------------------------------------------------

\paragraph{Descrizione}
L'attore \textit{Utente} rappresenta qualunque individuo che accede all'applicazione \productname{} prima di aver completato il processo di autenticazione. Questo profilo corrisponde allo stato iniziale di interazione con il sistema, limitato alla sola schermata di login.

\paragraph{Obiettivi}
\begin{itemize}[leftmargin=1.1em]
\item Autenticarsi al sistema utilizzando credenziali valide (email e password)
\item Accedere alle funzionalità della piattaforma secondo il proprio ruolo
\end{itemize}

\paragraph{Responsabilità e permessi}
\begin{itemize}[leftmargin=1.1em]
\item \textbf{Visualizzare} la schermata di autenticazione
\item \textbf{Inserire} email e password nei campi del form di login
\item \textbf{Inviare} una richiesta di autenticazione al sistema
\item \textbf{Ricevere} messaggi di errore in caso di credenziali non valide o account bloccato
\end{itemize}

\textbf{Limitazioni}: questo profilo non ha accesso ad alcuna funzionalità operativa del sistema fino al completamento dell'autenticazione.

\newpage

% ---------------------------------------------------------------------
\subsubsection{Utente\_Autenticato}
\label{subsec:utente-autenticato}
% ---------------------------------------------------------------------

\paragraph{Descrizione}

L'\textbf{Utente Autenticato} è un membro del team che ha effettuato correttamente il login e ha accesso ai \textbf{progetti a cui è stato assegnato da un amministratore}. Questo profilo rappresenta lo sviluppatore, tester o membro del team che lavora attivamente su issue all'interno dei progetti autorizzati.

\paragraph{Obiettivi}
\begin{itemize}[leftmargin=1.1em]
\item Visualizzare i progetti a cui ha accesso assegnato dall'amministratore
\item Lavorare sulle issue dei progetti autorizzati (creazione, modifica, commenti)
\item Collaborare con il team attraverso commenti e aggiornamenti sulle issue
\item Monitorare lo stato delle issue assegnate personalmente
\item Ricevere e gestire notifiche relative ai progetti di appartenenza
\end{itemize}

\paragraph{Responsabilità e permessi}
L'Utente\_Autenticato dispone dei seguenti permessi:

\subparagraph{Gestione Progetti}
\begin{itemize}[leftmargin=1.1em]
\item \textbf{Può}: visualizzare solo i progetti a cui è stato aggiunto da un amministratore, selezionare un progetto per accedere alle sue issue
\item \textbf{Non può}: creare nuovi progetti, eliminare progetti, modificare progetti, scegliere autonomamente a quali progetti partecipare, vedere progetti di cui non è membro
\end{itemize}

\subparagraph{Gestione Issue}
\begin{itemize}[leftmargin=1.1em]
\item \textbf{Creare} nuove issue, all'interno dei progetti a cui appartiene, specificando:
  \begin{itemize}
  \item Titolo (\textit{obbligatorio})
  \item Descrizione (\textit{obbligatoria})
  \item Tipo (\textit{obbligatorio}: Bug, Feature, Question, Documentation)
  \item Priorità (\textit{opzionale}: Nessuna (Default), Bassa, Media, Alta, Critica)
  \item Etichette personalizzate (\textit{opzionale}, array di stringhe)
  \item Immagini allegate (\textit{opzionale})
  \end{itemize}
\item \textbf{Visualizzare} la lista completa delle issue del progetto selezionato
\item \textbf{Filtrare} le issue per tipo, stato, priorità, etichetta, assegnatario, intervallo di data creazione
\item \textbf{Ordinare} le issue per data di creazione, priorità, stato (crescente/decrescente)
\item \textbf{Visualizzare} il dettaglio completo di una singola issue con tutti i suoi attributi
\item \textbf{Modificare} le issue esistenti (aggiornamento di titolo, descrizione, tipo, priorità, etichette), solo per i progetti di appartenenza
\item \textbf{Associare} etichette personalizzate a una issue per facilitarne la categorizzazione
\end{itemize}

\subparagraph{Collaborazione}
\begin{itemize}[leftmargin=1.1em]
\item \textbf{Aggiungere} commenti testuali a qualunque issue
\item \textbf{Eliminare} i propri commenti
\item \textbf{Leggere} i commenti lasciati da altri utenti
\end{itemize}

\paragraph{Vista Personalizzata}
\begin{itemize}[leftmargin=1.1em]
\item \textbf{Visualizzare} una vista dedicata delle issue assegnate al proprio account
\item \textbf{Ricevere} notifiche quando una issue viene assegnata al proprio account
\item \textbf{Visualizzare} l'elenco delle notifiche ricevute
\item \textbf{Confermare} la presa visione delle notifiche
\end{itemize}

\textbf{Limitazioni}: questo profilo non può creare o modificare utenze, non può assegnare issue ad altri utenti, non può archiviare issue e non ha accesso ai report amministrativi.

\newpage

% ---------------------------------------------------------------------
\subsection{Amministratore}
\label{subsec:amministratore}
% ---------------------------------------------------------------------

\subsubsection{Descrizione}
L'\textit{Amministratore} è un utente con privilegi elevati, responsabile della gestione delle utenze, della supervisione del flusso di lavoro e del monitoraggio delle performance del team. Questo ruolo eredita completamente tutti i permessi dell'Utente\_Autenticato e dispone di funzionalità amministrative aggiuntive.

\subsubsection{Obiettivi}
\begin{itemize}[leftmargin=1.1em]
\item \textbf{Gestire} gli accessi al sistema creando, modificando e rimuovendo account utente
\item \textbf{Garantire} la presa in carico tempestiva delle issue critiche attraverso l'assegnazione esplicita ai membri del team
\item \textbf{Supervisionare} l'andamento del progetto analizzando report mensili con metriche quantitative
\item \textbf{Mantenere} ordinate le viste principali archiviando le issue non più rilevanti ma preservandone la tracciabilità
\item \textbf{Monitorare} le performance individuali e del team per supportare decisioni strategiche
\end{itemize}

\subsubsection{Responsabilità e permessi}
L'Amministratore dispone di:

\paragraph{Permessi ereditati}
\begin{itemize}[leftmargin=1.1em]
\item \textbf{Tutti} i permessi e le responsabilità dell'Utente\_Autenticato (creazione issue, visualizzazione, filtri, commenti, notifiche)
\end{itemize}

\newpage

\paragraph{Gestione Utenze}
\begin{itemize}[leftmargin=1.1em]
\item \textbf{Creare} nuove utenze specificando:
  \begin{itemize}
  \item Nome (obbligatorio)
  \item Cognome (obbligatorio)
  \item Email (obbligatoria, univoca nel sistema)
  \item Password (obbligatoria, soggetta a policy di complessità)
  \item Ruolo (obbligatorio: Utente o Amministratore)
  \end{itemize}
\item \textbf{Modificare} i dati anagrafici e le credenziali di utenze esistenti
\item \textbf{Eliminare} utenze dal sistema previa conferma con inserimento della propria password amministrativa
\end{itemize}

\paragraph{Gestione Avanzata Issue}
\begin{itemize}[leftmargin=1.1em]
\item \textbf{Assegnare} una issue a un membro specifico del team, rendendolo responsabile della sua gestione
\item \textbf{Archiviare} una issue, spostandola dalle liste principali a una sezione di archivio dedicata (la issue rimane accessibile ma non appare nelle viste standard)
\end{itemize}

\paragraph{Analisi e Reportistica}
\begin{itemize}[leftmargin=1.1em]
\item \textbf{Visualizzare} report mensili contenenti metriche aggregate sull'attività del team:
  \begin{itemize}
  \item Numero totale di issue aperte nel periodo
  \item Numero totale di issue gestite/risolte nel periodo
  \item Tempo medio di risoluzione (aggregato e per singolo utente)
  \item Distribuzione delle issue per tipo e priorità
  \item Performance individuali di ciascun membro del team
  \end{itemize}
\end{itemize}

\textbf{Nota}: il sistema è fornito con un account amministratore predefinito con credenziali di default, che permette di effettuare il primo accesso e creare ulteriori utenze.

\newpage

% ---------------------------------------------------------------------
\subsection{Riepilogo Permessi per Profilo}
\label{subsec:riepilogo-permessi}
% ---------------------------------------------------------------------

La seguente tabella riassume i permessi differenziati per ciascun profilo utente:

\begin{longtable}{|p{0.35\linewidth}|c|c|c|}
\hline
\textbf{Funzionalità} & \textbf{Utente} & \textbf{Utente Aut.} & \textbf{Admin} \\ \hline
\endfirsthead

\hline
\textbf{Funzionalità} & \textbf{Utente} & \textbf{Utente Aut.} & \textbf{Admin} \\ \hline
\endhead

\hline
\endfoot

Autenticazione (login) & \cmark & \cmark & \cmark \\ \hline
Visualizza dashboard progetti & \xmark & \cmark (solo assegnati) & \cmark (tutti) \\ \hline
Crea progetto & \xmark & \xmark & \cmark \\ \hline
Modifica progetto & \xmark & \xmark & \cmark \\ \hline
Elimina progetto & \xmark & \xmark & \cmark \\ \hline
Gestisci membri progetto & \xmark & \xmark & \cmark \\ \hline
Seleziona progetto & \xmark & \cmark (solo assegnati) & \cmark (tutti) \\ \hline
Crea issue & \xmark & \cmark (progetti autorizzati) & \cmark \\ \hline
Visualizza issue & \xmark & \cmark (progetti autorizzati) & \cmark (tutti) \\ \hline
Modifica issue & \xmark & \cmark (progetti autorizzati) & \cmark (tutti) \\ \hline
Assegna issue & \xmark & \xmark & \cmark \\ \hline
Archivia issue & \xmark & \xmark & \cmark \\ \hline
Aggiungi commento & \xmark & \cmark & \cmark \\ \hline
Elimina proprio commento & \xmark & \cmark & \cmark \\ \hline
Filtra/ordina issue & \xmark & \cmark & \cmark \\ \hline
Aggiungi etichette/immagini & \xmark & \cmark & \cmark \\ \hline
Visualizza notifiche & \xmark & \cmark & \cmark \\ \hline
Crea utenza & \xmark & \xmark & \cmark \\ \hline
Modifica utenza & \xmark & \xmark & \cmark \\ \hline
Elimina utenza & \xmark & \xmark & \cmark \\ \hline
Visualizza report & \xmark & \xmark & \cmark \\ \hline

\end{longtable}

\textbf{Legenda}:
\begin{itemize}[leftmargin=1.1em]
\item \cmark{} = Permesso concesso
\item \xmark{} = Permesso negato
\item "solo assegnati" = può accedere solo a progetti in cui appare in \texttt{Project\_Member}
\item "progetti autorizzati" = operazioni limitate ai progetti di appartenenza
\end{itemize}

\newpage